\section{Paying every exceptional event at the oracle scale}
\label{sec:probability}

The preceding sections prove fitted-score comparisons on events whose
deterministic premises are consequences of raw noise, degree and row
concentration bounds. This section completes the unconditional theorems.
It is not enough that those events have probability tending to one.

\subsection{Events with an arbitrarily large fixed exponential rate}

For any fixed $H>0$, choose the raw norm threshold $L_H\sqrt d$ in
\eqref{eq:raw-norm-tail}. Since $\lambda_i\le\lambda_+d$ for a fixed
$\lambda_+$, the elementary Bernoulli Chernoff inequality also gives,
for a sufficiently large fixed $b_H$,
\[
 \PP\{\max_iD_i>b_Hd\}\le ne^{-Hd}.
\]
For example the exponent in the Poisson upper bound is
$b_H\log(b_H/\lambda_+)-b_H+\lambda_+$, which grows without bound as
$b_H$ grows. Here this is a bound on a Bernoulli transform, not a change
of sampling model.
The conditional events in \cref{lem:normalized-rows} contribute at most
$2Kn e^{-Hd}$. The constants in the deterministic lemmas may increase
with $H$, but remain independent of $n$.

For DC-SBM we additionally impose
\begin{equation}\label{eq:raw-block-event}
 \max_{a,b}|[Z^\top(A-\Omega)Z]_{ab}|\le L_Bn.
\end{equation}
Each centered block total is a sum of independent bounded upper-triangular
edge variables, with coefficients at most two and variance at most $Cnd$.
The completed diagonal contributes only $O(d)$.
Scalar Bernstein therefore implies, for sufficiently large fixed $L_B$
and all large $n$,
\[
 \PP\{\eqref{eq:raw-block-event}\text{ fails}\}
 \le 2K^2\exp\left\{-\frac{cL_B^2 n}{d+L_B}\right\}.
\]
Under $d^2\le C_1n$, the exponent is at least
$c'L_B^2d/C_1$ once $d\ge L_B$. Thus $L_B$ can be chosen to make this
probability at most $2K^2e^{-Hd}$ for any prescribed fixed $H$.
Ordinary free-profile fitting does not need this event: it integrates the
block-fluctuation exponential moment instead, as proved in
\cref{sec:sbm-proof}.

Let $\mathcal G$ be the intersection of these required events and the
minimum-degree event $D_i\ge a_0d$. All spectral, initialization and fitting
premises in the earlier sections follow on $\mathcal G$.
In particular all tested deleted graphs satisfy their own measurable
training premises there. Apart from the minimum-degree event, the failure
probability is at most
\begin{equation}\label{eq:other-bad-events}
 C_Kne^{-Hd}.
\end{equation}

\subsection{Ordinary SBM on the explicit parameter class}

For $P_{ab}\in[p/2,p]$ and $p\le1/4$, the log-odds coefficient of any
pairwise oracle contrast satisfies
\[
 |w_c^{ab}|\le W:=\log(7/3).
\]
Every tilted Bernoulli probability between the two candidates is at most
$p$, so the second derivative of the log affinity in
\eqref{eq:sbm-information} is at most $dW^2$.
A twice differentiable convex function vanishing at zero and one with
second derivative at most $M$ is at least $-Mt(1-t)/2$.
Consequently
\begin{equation}\label{eq:sbm-information-upper}
 I_{i,ab}\le d\,\frac{\log^2(7/3)}8<0.09d.
\end{equation}

For all sufficiently large $n$, every mean degree is at least $0.49d$.
Choose $a_0=1/16$. The Bernoulli lower-tail transform yields
\begin{align*}
 \PP(D_i<d/16)
 &\le \exp\left\{-d\left[
 0.49-\frac1{16}-\frac1{16}\log(16\cdot0.49)\right]\right\}\\
 &\le e^{-0.298d}.
\end{align*}
\begin{equation}\label{eq:sbm-low-degree}
 \PP(\min_iD_i<d/16)\le ne^{-0.298d}.
\end{equation}
Take $H=1$, increasing the other fixed thresholds if necessary.
Equations \eqref{eq:other-bad-events} and
\eqref{eq:sbm-low-degree} give
$\PP(\mathcal G^c)\le Cne^{-0.298d}$.
The information bound \eqref{eq:sbm-information-upper} applies to every
vertex and every competitor. Thus the sharp pairwise lower bound
\eqref{eq:oracle-pairs}, followed by averaging over all vertices, gives
\[
 R_{\mathrm{SBM},*}\ge
 \frac1n\sum_{i=1}^n\frac c{\sqrt d}e^{-0.09d}
 =\frac c{\sqrt d}e^{-0.09d}.
\]
No additional balancing assumption is needed: the uniform coefficient
bound, group-proportion bound and full-rank separation already hold at
every vertex. Consequently
\[
 \frac{\PP(\mathcal G^c)}{R_{\mathrm{SBM},*}}
 \le Cn\sqrt d\,e^{-0.208d}.
\]
The function $\sqrt d\,e^{-0.208d}$ is decreasing once
$d>1/(2\cdot0.208)$. Therefore, for $d\ge5\log n$ and all sufficiently
large $n$,
\[
 n\sqrt d\,e^{-0.208d}
 \le \sqrt{5\log n}\,n^{1-5\cdot0.208}
 =\sqrt{5\log n}\,n^{-0.04}=o(1).
\]
The earlier coefficient $20$ was a sufficient numerical choice obtained
by retaining only one vertex in the oracle lower bound and then bounding
$\sqrt d$ by $\sqrt n$. Both losses are unnecessary. The improved
coefficient $5$ changes neither the encoder nor the likelihood decoder;
it does not remove the global minimum-degree event from this proof.
The good-event comparison in \cref{sec:sbm-proof} and
\eqref{eq:bad-payment} complete \cref{thm:sbm}.

\subsection{The automatic margin for a general DC profile oracle}

Let $J_{i,a\to b}$ be the exact Bernoulli directional profile information
defined in \cref{sec:dcsbm-proof}. For the self-excluded neighbor masses
$M_c=\sum_{j\ne i:z_j=c}\theta_j$, set
\[
 T_a=\sum_cM_cC_{ac},\qquad
 q_{ac}=\frac{M_cC_{ac}}{T_a},\qquad
 h_{ab}^*=\min_{0\le s\le1}\sum_cq_{ac}^{1-s}q_{bc}^s.
\]
Uniform positive group proportions and entries of $C$ give
$q_{ac}\ge q_->0$ for all large $n$. Hence
\begin{equation}\label{eq:profile-overlap}
 h_{ab}^*\ge\sum_c\min(q_{ac},q_{bc})\ge Kq_->0.
\end{equation}
The exact log-transform expansion, uniformly in $s\in[0,1]$, is
\begin{align*}
 \sum_{j\ne i}\log\{1+\theta_i\theta_jpC_{az_j}
                    (e^{-sw_{z_j}}-1)\}
 &=\theta_ipT_a
       \left[\sum_cq_{ac}^{1-s}q_{bc}^s-1\right]+O(np^2).
\end{align*}
Minimizing therefore gives
\begin{equation}\label{eq:profile-leading-information}
 J_{i,a\to b}
 =\lambda_i(1-h_{ab}^*)+O(np^2),
 \qquad \lambda_i=\theta_ipT_a.
\end{equation}
This calculation uses the exact Bernoulli transform.
Under $np^2\le C_1$, \eqref{eq:profile-overlap} and
$\lambda_i\ge\lambda_-d$ imply, for a fixed $\eta>0$ and all large $n$,
\begin{equation}\label{eq:profile-degree-margin}
 J_{\min}:=\min_i\min_{b\ne z_i}J_{i,z_i\to b}
 \le\lambda_{\min}-\eta d.
\end{equation}
To see it, choose a vertex with minimum mean degree and any competitor.
Equal candidate exposures are not required for this profile margin.

Choose $a_0>0$ sufficiently small that
\[
 a_0\{1+\log(\lambda_+/a_0)\}\le\eta/2,\qquad a_0<\lambda_-.
\]
For each vertex, the Bernoulli lower-tail bound is
\[
 \PP(D_i<a_0d)
 \le\exp\{-\lambda_i+a_0d+
                   a_0d\log(\lambda_i/(a_0d))\}.
\]
The exponent is monotone in $\lambda_i$ above $a_0d$.
Together with \eqref{eq:profile-degree-margin}, it follows that
\begin{equation}\label{eq:dc-low-degree}
 \PP(\min_iD_i<a_0d)
 \le n\exp\{-J_{\min}-\eta d/2\}.
\end{equation}
The full profile reference has the lower bound
\[
 R_{\mathrm{prof},*}\ge
 \frac c{n\sqrt d}e^{-J_{\min}}.
\]
Thus the ratio of \eqref{eq:dc-low-degree} to that reference is at most
\[
 Cn^2\sqrt d\,e^{-\eta d/2}=o(1)
\]
when $d\ge C_0\log n$ with, for example, $C_0>6/\eta$ and sufficiently
large $n$. We used only $d\le n$ in bounding the polynomial factor.
The constants can be enlarged without changing their independence of $n$.

Also $J_{\min}\le C_Jd$ for a fixed $C_J$ by its bounded-weight transform.
Choosing $H>C_J+1$ in \eqref{eq:other-bad-events}, and enlarging $C_0$ if
needed, pays those remaining failures at the same reference scale.
Hence $\PP(\mathcal G^c)=O(R_{\mathrm{prof},*})$.
The good-event theorem and invariant-basin result in
\cref{sec:dcsbm-proof} now prove \cref{thm:dc}, including its finite
hard-refit version.

Under the equal-exposure condition, the profile and known-target-degree
oracle risks are comparable by the bounded-contrast proof in
\cref{sec:dcsbm-proof}. This proves the DC assertion of \cref{cor:known}.
For the native Euclidean ordinary-SBM decoder, the same sparse transform
expansion gives
$I_{i,ab}=\lambda_i(1-h_{ab}^*)+O(np^2+p)$ under equal exposures.
Thus its information--degree margin is also automatic.
The general-positive-profile version of \cref{sec:sbm-proof}, with its
block fluctuation integrated instead of imposing \eqref{eq:raw-block-event},
therefore proves the remaining assertion of \cref{cor:known}.

\subsection{An explicit extension for ordinary SBM with unequal exposures}

The entry-ratio-two restriction in \cref{thm:sbm} is a simple sufficient
way to ensure the failure-event comparison. The deterministic and
conditional likelihood proof covers a larger parameter range.

\begin{proposition}[A deterministic information--degree margin suffices]
\label{prop:sbm-margin}
Suppose ordinary SBM has fixed $K$, positive comparable group sizes,
$c_-p\le P_{ab}\le c_+p$, $\max P_{ab}\le1-\epsilon_0$,
and $\sigma_{\min}(P)\ge\kappa p$, for fixed positive constants.
If for some fixed $\eta>0$ and all sufficiently large $n$,
\[
 \lambda_{\min}-I_{\min}\ge\eta d,\qquad d=np\ge C_0\log n,
\]
then the native Euclidean signed likelihood decoder, with sufficiently
small structural floor coefficient, satisfies
$\E\ell(\widehat z_E,z)\le C R_{\mathrm{SBM},*}$
when $C_0$ is sufficiently large. No condition $np^2=O(1)$ is needed
for this ordinary-SBM proposition.
\end{proposition}

\begin{proof}
All true Bernoulli likelihood coefficients and tilted means are exactly
block constant, so the Euclidean projection, fitted-score and sharp-tail
arguments of \cref{sec:sbm-proof} apply without a sparse Bernoulli
approximation. Generated profiles remain comparable to $p$ and uniformly
away from one when the floor is chosen for the structural class.
The displayed margin, the same small degree cutoff used above and an
arbitrarily strong choice of the other exponential bounds imply
$\PP(\mathcal G^c)=O(R_{\mathrm{SBM},*})$.
Here the polynomial comparison uses $d\le Cn$, which follows from
$c_-p\le P_{ab}\le1-\epsilon_0$; it does not require $p\le1$.
\end{proof}

This proposition is a condition on the deterministic model parameters,
not an assumed success event or a hypothesis that the desired risk ratio
already holds. It is nevertheless an additional restriction. For model
families violating it, a refined analysis of exceptional low-degree
graphs or a different normalized encoder is needed; the crude loss bound
one on all such graphs does not prove a constant ratio.
